\documentclass[11pt,letterpaper]{article}

% --- Packages ---
\usepackage[margin=1in]{geometry}
\usepackage{booktabs}
\usepackage[colorlinks=true,linkcolor=blue,urlcolor=blue,citecolor=blue]{hyperref}
\usepackage{amsmath}
\usepackage{listings}
\usepackage{xcolor}
\usepackage{graphicx}
\usepackage{siunitx}
\usepackage{enumitem}
\usepackage{longtable}
\usepackage{caption}

% --- Listings style ---
\lstset{
  basicstyle=\ttfamily\small,
  backgroundcolor=\color{gray!10},
  frame=single,
  framerule=0pt,
  breaklines=true,
  columns=fullflexible,
  keepspaces=true,
  aboveskip=1em,
  belowskip=1em,
}

% --- siunitx setup ---
\sisetup{
  per-mode=symbol,
  detect-all,
}

% --- Title ---
\title{Signal-to-Noise Improvement Recommendations}
\author{LDRD FFA@CEBAF Permanent Magnet Radiation Damage Study}
\date{June 2026}

\begin{document}
\maketitle

\noindent\textbf{Audience:} Summer interns and team members working on
measurement improvements.

\bigskip\noindent\textbf{Figures:} This document references PNG files from the
project. For Overleaf, upload them to the same directory as this \texttt{.tex}
file. Source locations relative to \texttt{Claude\_Analysis1/}:
\begin{itemize}[nosep]
  \item \texttt{Measurement\_Improvement/} (TC-series)
  \item \texttt{Cleanup\_Claude/Presentation\_Plots/} (P-series)
\end{itemize}

\tableofcontents
\newpage

%=============================================================================
\section{Current Noise Sources (Ranked by Impact)}
\label{sec:noise}
%=============================================================================

\subsection{Teslameter Positioning Noise (0.14--1.90\%)}
\label{sec:noise-positioning}

\textbf{Impact:} Dominant noise source for Teslameter field measurements.

The 3D-printed sample rigs allow ${\sim}$0.1--\SI{0.3}{\milli\meter} of play
when a magnet is inserted. Because the Hall probe measures field at a single
point, this positional uncertainty translates directly to measurement scatter:

\begin{table}[htbp]
\centering
\caption{Teslameter positioning noise by face.}
\label{tab:positioning-noise}
\begin{tabular}{@{}lccc@{}}
\toprule
Face & Std (\%) & Typical reading (mT) & $1\sigma$ noise (mT) \\
\midrule
Top   & 0.14--0.29 & ${\sim}$150--180 & ${\sim}$0.2--0.5 \\
Front & 0.57--0.91 & ${\sim}$30--50   & ${\sim}$0.2--0.5 \\
Side  & 0.90--1.90 & ${\sim}$10--30   & ${\sim}$0.1--0.6 \\
\bottomrule
\end{tabular}
\end{table}

The top face has the flattest field gradient, so positioning errors matter least
there. Front and side faces sit in steeper gradient regions.

\textbf{Current status:} This noise does not affect Helmholtz demagnetization
results (the Helmholtz integrates over the full sample volume). However, the
Teslameter is NOT merely a secondary cross-check: if positioning noise can be
controlled, Teslameter data becomes an independent primary measurement of
spatially-resolved point-field degradation. Showing both integrated moment loss
(Helmholtz) and point-field loss (Teslameter) on the same samples, across
Y-plates, H-plates, and A-samples, is the path to high-impact publications.

\textbf{Why this matters for H-plates especially:} H-plates are Halbach-geometry
assemblies that directly represent FFA magnet configurations. The Teslameter
measures field at specific locations on the Halbach assembly; demonstrating
degradation in the actual field geometry of interest is far more compelling to
accelerator designers than volume-averaged moment loss alone.

\textbf{Quantitative target:} Reduce top-face noise to $<0.15\%$, which would
make the Teslameter precise enough to independently detect the ${\sim}0.2\%$
demagnetization signal. At $<0.10\%$, it would be comparable to Helmholtz
single-session repeatability.


\subsection{Temperature Uncertainty (up to \SI{3.4}{\degreeCelsius} Sensor Disagreement)}
\label{sec:noise-temp}

\textbf{Impact:} $0.066\%/\si{\degreeCelsius}$ on the NdFeB--SmCo differential;
up to 0.033\% at current $\pm$\SI{0.5}{\degreeCelsius} uncertainty level.

June 2--3, 2026 data (492 matched pairs from
\texttt{temp\_sensor\_comparison.py}) shows a systematic offset between the
Teslameter probe and Arduino AM2302 readings:

Fleet statistics (Teslameter minus Arduino):
\begin{itemize}[nosep]
  \item Mean: \textbf{+\SI{1.00}{\degreeCelsius}}
  \item Std: \SI{0.74}{\degreeCelsius}
  \item Range: $-1.01$ to $+$\SI{3.43}{\degreeCelsius}
\end{itemize}

Temperature-dependent pattern:
\begin{itemize}[nosep]
  \item Low temp ($<$\SI{24}{\degreeCelsius}, $N=17$): mean offset $= +$\SI{1.87}{\degreeCelsius}
  \item High temp ($\geq$\SI{27}{\degreeCelsius}, $N=401$): mean offset $= +$\SI{0.86}{\degreeCelsius}
  \item The offset shrinks as both sensors converge at higher temperatures.
\end{itemize}

Figures~\ref{fig:tc1} and~\ref{fig:tc3} illustrate these patterns.

\begin{figure}[htbp]
  \centering
  \includegraphics[width=0.85\textwidth]{TC1_temp_scatter.png}
  \caption{Teslameter vs Arduino temperature (492 matched pairs, June 2--3,
    2026). The Teslameter reads systematically higher, especially at low
    temperatures.}
  \label{fig:tc1}
\end{figure}

\begin{figure}[htbp]
  \centering
  \includegraphics[width=0.85\textwidth]{TC3_timeofday.png}
  \caption{Temperature offset (Teslameter minus Arduino) vs time of day. The
    disagreement is largest in the early morning and shrinks through the day.}
  \label{fig:tc3}
\end{figure}

\clearpage
Y-14 calibration plate examples:

\begin{table}[htbp]
\centering
\caption{Y-14 temperature readings: Arduino vs Teslameter.}
\label{tab:y14-temp}
\begin{tabular}{@{}lccc@{}}
\toprule
Measurement & Arduino (\si{\degreeCelsius}) & Teslameter (\si{\degreeCelsius}) & Diff (\si{\degreeCelsius}) \\
\midrule
Jun 1 10:43 & 21.1 & 21.6 & $+0.5$ \\
Jun 2 09:04 & 20.3 & 22.7 & $+2.4$ \\
Jun 3 07:27 & 22.6 & 24.0 & $+1.4$ \\
Jun 3 08:32 & 27.4 & 27.9 & $+0.5$ \\
Jun 3 09:25 & 28.5 & 28.3 & $-0.2$ \\
\bottomrule
\end{tabular}
\end{table}

\textbf{Multiple possible explanations} (not yet resolved):
\begin{itemize}[nosep]
  \item Probe self-heating (Teslameter biased high)
  \item Thermal lag: samples may genuinely be warmer than the air early in the
    day (residual heat from previous day, slow equilibration), meaning the
    Teslameter probe, being in contact with the sample, is closer to the true
    sample temperature
  \item Arduino sensitivity to local air currents (cross-breezes, laptop fans,
    HVAC) could bias the Arduino reading away from the actual temperature near
    the sample
\end{itemize}

\textbf{Which sensor is ``right'' for the temperature correction is an open
question.} The correction formula needs the \emph{sample} temperature, not
necessarily the air temperature. The Teslameter probe, being in thermal contact
with the sample surface, may actually be the better proxy in some conditions. The
summer cross-calibration study is designed to disentangle these effects.

Linear fit: $T_{\mathrm{tesla}} = 0.707 \times T_{\mathrm{arduino}} + 9.05$
(fit residual std $= \SI{0.62}{\degreeCelsius}$).

\textbf{Quantitative target:} Understand the source of the disagreement well
enough to select (or combine) sensors with confidence. Target: reduce
temperature uncertainty to $<$\SI{0.3}{\degreeCelsius}, which would bring the
temperature systematic on the differential from $\pm 0.033\%$ to $\pm 0.020\%$.


\subsection{Session-to-Session Gain Drift (0.5--0.7\%)}
\label{sec:noise-gain}

\textbf{Impact:} $\pm 0.124\%$ systematic on individual material measurements.
Cancels to zero in the intra-plate differential.

Measured from 5 pre-deployment sessions. The drift is a multiplicative shift in
the Helmholtz system output (coil parameter changes, cable connections,
electronics temperature). Nov 2024 sessions read ${\sim}0.7\%$ higher than
April--June 2025 sessions.

\textbf{Current status:} Fully handled by the intra-plate differential for
Y-plates. NOT handled for H-plates (no intra-plate differential possible) or
absolute measurements.

\textbf{Quantitative target:} Track and correct gain within sessions using a
reference magnet, reducing individual-material systematic to $<0.05\%$.


\subsection{First-Tunnel Bias for H-Plates (62\% of Front Readings $>$2\% Low)}
\label{sec:noise-htunnel}

\textbf{Impact:} Corrupts the first tunnel measurement for H-plate front and
side faces, forcing use of only the second and subsequent measurements.

The first time H-plates were measured in the tunnel (July 17, 2025), a
systematic low bias appeared on front (62\%) and side (61\%) faces. Top face was
unaffected. Likely caused by different rig setup or insufficient equilibration on
the first trip.

\textbf{Current status:} Mitigated by using only post-first-tunnel data for
H-plates. Reduced effective baseline for some H-plates.

\textbf{Quantitative target:} Identify root cause (rig setup procedure? thermal
equilibration?) and eliminate for C2.


\subsection{Single-Session Baselines}
\label{sec:noise-baseline}

\textbf{Impact:} Most tunnel plates have only 1 qualifying pre-deployment
session with co-located Teslameter temperature. Baseline noise is not averaged
down.

With a single baseline, the baseline uncertainty is the single-session
repeatability (${\sim}$0.2--0.3\% for Helmholtz). With $N=5$ independent
sessions, this drops to ${\sim}$0.1--0.13\%. The improvement is roughly
$1/\sqrt{N}$.

Example: Y-10 had 2 qualifying baseline sessions. Its per-material scatter
dropped by a factor of ${\sim}$20 compared to single-baseline plates. (This is
partially because the two baselines happened to average out a measurement fluke,
but it demonstrates the principle.)

\textbf{Quantitative target:} 5--10 independent baseline sessions per lab plate.


%=============================================================================
\section{Recommended Improvements}
\label{sec:improvements}
%=============================================================================

\subsection{Multi-Insertion Averaging (Teslameter)}
\label{sec:imp-multiinsert}

\textbf{What:} For each Teslameter measurement, remove the sample from the rig,
re-insert it, and re-measure. Repeat $3\times$ per face. Report the mean
$\pm$~std.

\textbf{Why:} Directly measures the rig positioning tolerance for each
individual sample, in real time. The mean of 3 insertions has
$\sqrt{3} = 1.7\times$ better precision than a single insertion. More
importantly, the per-sample std tells you whether that specific sample has good
or poor rig fit.

\textbf{Protocol:}
\begin{enumerate}[nosep]
  \item Insert sample into rig, measure top face ($B_x$, $B_y$, $B_z$, $T$)
  \item Remove sample completely from rig
  \item Re-insert sample, measure top face again
  \item Repeat once more (3 total insertions)
  \item Record all 3 readings and compute mean, std
  \item Repeat for front and side if desired
\end{enumerate}

\textbf{Expected improvement:} Top-face std from 0.14--0.29\% (single insertion)
to 0.08--0.17\% (mean of 3). Samples with std $>$0.3\% can be flagged for rig
inspection.

\textbf{Time cost:} ${\sim}$2 extra minutes per sample per face. For 120
Y-plate slots, this adds ${\sim}$4 hours to a full Teslameter survey if doing
top face only.


\subsection{Top-Face Priority}
\label{sec:imp-topface}

\textbf{What:} For routine trend tracking, measure only the top face on the
Teslameter. Reserve front and side for dedicated characterization campaigns.

\textbf{Why:} Top face is 3--6$\times$ more precise than other faces. For
tracking demagnetization trends over time, top-face-only provides adequate
precision with ${\sim}1/3$ the measurement time.

\textbf{When to do all 3 faces:} Initial characterization, dedicated
repeatability studies, and whenever a result looks anomalous on top face.


\subsection{Dual Temperature Sensors}
\label{sec:imp-dualtemp}

\textbf{What:} Record both Arduino AM2302 and Teslameter temperatures for every
measurement. Neither is designated as ``primary'' until the cross-calibration
study determines which (or what combination) best represents the sample
temperature for the correction formula.

\textbf{Why:} The two sensors measure different things:
\begin{itemize}[nosep]
  \item \textbf{Arduino AM2302:} Ambient air temperature near the Helmholtz
    setup. May be influenced by cross-breezes, laptop fans, HVAC airflow, and
    distance from the sample.
  \item \textbf{Teslameter probe:} Temperature at the probe junction, in contact
    with (or very near) the magnet surface. May include self-heating bias, but is
    physically closer to the sample whose temperature we actually need.
\end{itemize}

Having both sensors lets us: (a)~quantify their disagreement, (b)~study which
correlates better with Helmholtz reading residuals, and (c)~eventually select or
combine them for optimal temperature correction.

\textbf{Protocol:}
\begin{enumerate}[nosep]
  \item Arduino sensor runs continuously, logging to
    \texttt{*\_helmholtzTemp.dat} files
  \item Arduino reading taken at the moment of Helmholtz measurement
  \item Teslameter temperature recorded when sample is in Teslameter rig (may be
    minutes later)
  \item Both temperatures recorded in the data file for each sample
\end{enumerate}

\textbf{Cross-calibration script:}
\texttt{Measurement\_Improvement/temp\_sensor\_comparison.py} uses the June 2--3
data to characterize the Arduino--Teslameter offset pattern.

\textbf{Expected improvement:} Once the relationship is understood, the
temperature systematic could drop from $\pm 0.033\%$ toward $\pm 0.020\%$ on the
differential. But this depends on understanding which sensor better represents
the sample temperature, not simply assuming one is correct.


\subsection{Environmental Conditioning}
\label{sec:imp-enviro}

\textbf{What:} Let samples equilibrate to room temperature before measurement.
Record the wait time. Use Arduino humidity data to flag condensation risk.

\textbf{Why:} If a sample is cooler or warmer than the room, the Teslameter
temperature (which reads the magnet surface) will not match the Arduino
temperature (which reads the air). Equilibration eliminates this discrepancy.

\textbf{Protocol:}
\begin{enumerate}[nosep]
  \item Place samples on bench in measurement room at least 30 minutes before
    starting
  \item Record the time the samples were moved and the time measurement begins
  \item Note Arduino humidity: if $>$60\% and sample is cold, condensation risk
    exists (avoid measurement until humidity drops or sample warms)
  \item For high-precision measurements, wait 1 hour and verify
    Arduino--Teslameter agreement is within \SI{0.5}{\degreeCelsius}
\end{enumerate}


\subsection{Multi-Session Baselines}
\label{sec:imp-baselines}

\textbf{What:} Measure every lab plate in 5--10 independent sessions over 5
weeks. Use the mean as the baseline.

\textbf{Why:} Noise reduces as $1/\sqrt{N}$. With $N=10$ sessions, the baseline
noise contribution drops from ${\sim}$0.2--0.3\% (single session) to
${\sim}$0.06--0.10\%.

\textbf{Quantitative target:}
\begin{itemize}[nosep]
  \item Current lab SEM: 0.038\% (9 plates, mostly single-session baselines)
  \item Target lab SEM: ${\sim}$0.012\% (9 plates, 10-session baselines)
  \item This would make the lab null result definitive at the 0.01\% level
\end{itemize}

\textbf{Cadence:} 2 sessions per week for 5 weeks. Each session measures all 9
lab Y-plates (36 slots), all 48 lab H-plates (192 slots), and lab A-samples.
Arduino sensor on for all Helmholtz measurements.


\subsection{Rig Improvement Assessment}
\label{sec:imp-rig}

\textbf{What:} Directly measure rig backlash and positioning tolerance using a
reference magnet. If $>$0.3\%, consider designing a machined aluminum jig for
top-face measurements.

\textbf{Protocol:}
\begin{enumerate}[nosep]
  \item Select 1 reference magnet (e.g., a spare N42EH slug)
  \item Insert into rig, measure top face. Do NOT remove.
  \item Gently push sample to one extreme of rig play, measure again
  \item Push to opposite extreme, measure again
  \item Repeat 10$\times$, characterize the range of readings
\end{enumerate}

\textbf{Decision threshold:} If rig backlash alone accounts for $>$0.3\%
variation, a machined jig is justified. A simple aluminum V-groove or flat-plate
holder with set screw could reduce positioning noise to $<$0.05\%.


%=============================================================================
\section{Helmholtz Improvement Ideas}
\label{sec:helmholtz-improve}
%=============================================================================

\subsection{Multiple Insertions per Helmholtz Measurement}
\label{sec:helm-multiinsert}

\textbf{What:} Remove sample from Helmholtz coil, re-insert, re-measure. Repeat
3--5 times per sample.

\textbf{Why:} Helmholtz readings are position-independent ``by design,'' but
small variations in sample centering may still contribute. Multiple insertions
quantify the single-insertion reproducibility directly.

\textbf{Expected result:} We expect the single-insertion std to be 0.02--0.05\%
(much better than Teslameter), confirming that Helmholtz positioning noise is
negligible compared to gain drift.

\textbf{Time cost:} ${\sim}$1 extra minute per insertion. For 120 samples
$\times$ 3 extra insertions $= {\sim}$6 hours additional. Best done with a
subset first (Y-14 + 2 other lab plates, 12 slots total, ${\sim}$36 extra
insertions in 30 minutes).


\subsection{Same-Day Reference Magnet}
\label{sec:helm-refmag}

\textbf{What:} Measure a designated reference magnet (always the same physical
sample) at the start, middle, and end of every session. Use it to track
within-session gain drift.

\textbf{Why:} Currently, within-session drift is an upper bound (we showed it is
consistent with zero, but did not directly measure it). A reference magnet gives
a direct, real-time gain monitor.

\textbf{Protocol:}
\begin{enumerate}[nosep]
  \item Designate one lab magnet as the session reference (e.g., Y-14 slot~1, or
    a standalone slug with known properties)
  \item Measure it first thing, after every 10 plates, and last thing
  \item Plot the reference readings vs time within the session
  \item If drift is detected, apply a linear correction to all other readings
\end{enumerate}

\textbf{Expected benefit:} Transforms the gain systematic from a session-wide
unknown into a measured, correctable quantity. Could reduce
individual-material systematic from $\pm 0.124\%$ to $<\!\pm 0.05\%$.


\subsection{Synchronous Temperature at Helmholtz Measurement}
\label{sec:helm-synctemp}

\textbf{What:} The Arduino AM2302 records temperature at the moment of each
Helmholtz measurement, while the Teslameter temperature is recorded 1--5 minutes
later (when the sample is moved to the Teslameter rig). Having a temperature
reading synchronous with the Helmholtz measurement is valuable regardless of
which sensor ultimately proves more accurate.

\textbf{Why:} The Teslameter reading is time-delayed relative to the Helmholtz
measurement, during which the room temperature may change. However, the Arduino
measures ambient air, not sample temperature; the Teslameter probe is in contact
with the sample. In a rapidly warming room, the Arduino may track the air faster
while the sample lags; in a stable room, the time delay matters less.

\textbf{Implementation:} The \texttt{*\_helmholtzTemp.dat} files already contain
per-sample Arduino temperatures. The analysis pipeline should be updated to
include both temperature sources as columns, and the cross-calibration study
should determine which (or what weighted combination) produces smaller Helmholtz
residuals.


%=============================================================================
\section{Measurement Priority Matrix}
\label{sec:priority-matrix}
%=============================================================================

\begin{table}[htbp]
\centering
\caption{Improvement priority matrix. ``Differential'' refers to the gain-immune
  Y-plate headline result; ``Individual'' refers to per-material tracking,
  H-plate analysis, and publication scope.}
\label{tab:priority}
\small
\begin{tabular}{@{}p{3.2cm}lp{2.5cm}p{2.5cm}c@{}}
\toprule
Improvement & Effort & Impact on differential & Impact on individual & Priority \\
\midrule
Multi-session baselines      & High (5 weeks)         & Moderate (stat only)      & Large (stat + gain)                                  & \textbf{1} \\
Dual temperature sensors     & Low (already in place)  & Moderate (if offset understood) & Moderate                                      & \textbf{2} \\
Multi-insertion Teslameter   & Medium (extra 4 hrs/survey) & None (Helmholtz-only) & \textbf{HIGH} (enables dual-measurement publications) & 3 \\
Multi-insertion Helmholtz    & Low (30 min pilot)      & Small (quantify floor)    & Small                                                & \textbf{4} \\
Reference magnet tracking    & Low--Medium              & None (diff.\ already 0)  & Large (gain control)                                 & \textbf{5} \\
Rig assessment               & Low (1 hour)            & None                      & \textbf{HIGH} (gates Teslameter as science)           & 6 \\
Environmental conditioning   & Low (protocol change)   & Small ($<$0.01\%)        & Small                                                & 7 \\
\bottomrule
\end{tabular}
\end{table}

\textbf{Note:} ``Impact on differential'' matters for the headline Y-plate
science result. ``Impact on individual'' matters for per-material tracking,
H-plate analysis, and the broader publication story. For high-tier publications,
the goal is to show consistent demagnetization signatures across BOTH measurement
techniques (integrated + point-field) and all three sample types (Y, H, A). The
Teslameter improvements are what unlock this.


%=============================================================================
\section{Expected Outcomes}
\label{sec:outcomes}
%=============================================================================

\subsection{With Full Implementation (End of Summer)}

\begin{table}[htbp]
\centering
\caption{Current vs target measurement parameters.}
\label{tab:outcomes}
\begin{tabular}{@{}lp{2.5cm}p{2.5cm}p{4.5cm}@{}}
\toprule
Parameter & Current & Target & Method \\
\midrule
Lab Y-plate differential SEM          & 0.038\%            & 0.012\%                       & 10-session baselines \\
Temperature systematic                 & $\pm 0.033\%$     & $\pm 0.020\%$                 & Dual-sensor cross-calibration \\
Helmholtz single-insertion std         & Unknown            & Measured (expect 0.02--0.05\%) & Multi-insertion study \\
Teslameter top-face std                & 0.14--0.29\%      & $<$0.15\% (measured)           & Multi-insertion + rig characterization \\
Individual material gain syst.         & $\pm 0.124\%$     & $\pm 0.05\%$                  & Reference magnet protocol \\
\bottomrule
\end{tabular}
\end{table}

\subsection{Impact on Science and Publications}

\textbf{Near-term (Tech Note, IPAC-class):} With reduced lab SEM, the lab null
result becomes definitive, strengthening the existing C1 conclusion.

\textbf{High-impact (journal publication):} The transformative outcome is if the
Teslameter noise reduction succeeds. Currently:

\begin{itemize}[nosep]
  \item We have a strong Helmholtz result (integrated moment, gain-immune
    differential, $7.6\sigma$) but only on Y-plates.
  \item H-plate Helmholtz data carries the full gain systematic ($\pm 0.124\%$)
    because H-plates have no intra-plate differential. The gain dominates and the
    results are noisy.
  \item Teslameter data is currently too noisy (0.14--1.90\%) to independently
    confirm the ${\sim}0.2\%$ signal.
\end{itemize}

If the summer improvements achieve $<$0.15\% Teslameter top-face precision AND
multi-session baselines bring the gain systematic under control:

\begin{enumerate}[nosep]
  \item \textbf{Dual-measurement confirmation:} The same demagnetization signal
    seen in both integrated moment AND spatially-resolved point field, on the
    same samples. This is a qualitatively different level of evidence.
  \item \textbf{H-plate Halbach results:} Degradation measured in the actual
    field geometry relevant to FFA magnets, using both techniques. This directly
    answers the design question that motivates the entire study.
  \item \textbf{A-sample paired comparisons:} Field-vs-zero-field degradation in
    identical material pairs. Combined with both measurement techniques, this
    constrains the degradation mechanism.
  \item \textbf{Per-face, per-axis spatial information:} Where on the magnet is
    the field lost? Is it uniform or concentrated? The Teslameter's 3-axis,
    3-face data provides this; the Helmholtz cannot.
\end{enumerate}

The difference between ``Y-plate Helmholtz differential'' and ``Y + H + A
samples, Helmholtz + Teslameter, with noise characterized and controlled'' is the
difference between a solid technical note and a definitive journal publication.


\end{document}
